\chapter[Introdução]{Introdução}
\label{cap:introducao}

O registro eletrônico em saúde concentra grande parte da informação clínica
sob a forma de texto livre. Notas de evolução, sumários de alta e laudos são
redigidos em linguagem natural, com abreviações idiossincráticas, ausência
frequente de pontuação padronizada e construções negativas características do
discurso médico. Essa riqueza textual, embora valiosa para pesquisa e apoio à
decisão, permanece majoritariamente inacessível a consultas estruturadas.

A extração de informações estruturadas desses textos é uma tarefa desafiadora, sobretudo em razão de suas características linguísticas, uma vez que um mesmo achado pode ser descrito de diferentes formas, apresentado de maneira abreviada ou expresso por meio de afirmações e negações em frases linguisticamente semelhantes. Técnicas de Processamento de Linguagem Natural (\foreignlanguage{english}{Natural Language Processing}, NLP) buscam transformar o texto clínico em conhecimento estruturado, mas, no
contexto da língua portuguesa, seu avanço foi historicamente limitado pela escassez
de recursos anotados. A disponibilização do corpus \textbf{SemClinBr}
\cite{oliveira_2022_semclinbr}, primeiro corpus clínico anotado semanticamente para o
português, com mil documentos provenientes de múltiplas instituições e
especialidades, criou condições para investigar tarefas antes restritas ao inglês,
entre elas a extração de relações entre entidades clínicas.

\section{Problema de pesquisa}

Uma primeira etapa da extração de informação clínica é o reconhecimento de entidades
nomeadas (\foreignlanguage{english}{Named Entity Recognition}, NER), que identifica
no texto as menções a conceitos clínicos (sintomas, doenças, achados,
procedimentos) e classifica cada uma em uma categoria. Em uma frase como ``paciente
sem febre'', o NER pode reconhecer ``febre'' como sinal ou sintoma e ``sem'' como uma
expressão de negação.

Identificar as entidades, porém, não basta para compreender o que o texto afirma
sobre elas. As duas menções do exemplo, isoladas, não informam que é a febre que está
sendo negada, nem distinguem essa frase de outra em que a febre é relatada e a
negação se refere a um conceito vizinho. Essa articulação entre entidades é o objeto
da extração de relações (\foreignlanguage{english}{Relation Extraction}, RE), etapa
complementar ao NER que determina se duas entidades já identificadas mantêm uma
relação e de que tipo.

A negação torna essa etapa especialmente relevante no domínio clínico. Não basta
encontrar a expressão que indica ausência ou negação: é preciso identificar a
informação clínica à qual ela se refere. Tratar como presente um achado que o texto
nega pode inverter a interpretação de um prontuário, e um mesmo documento costuma
conter várias expressões de negação, cada uma ligada a um conceito diferente.

Para o português, a extração de relações clínicas é pouco representada na
literatura, que se concentra na classificação de textos e no reconhecimento de
entidades \cite{shimaoka_2026_revisao}. O problema abordado neste trabalho consiste, portanto, no desenvolvimento de um modelo de extração de relações em notas clínicas em português, com ênfase na identificação de relações de negação. Esse modelo, denominado RECLin-PT, é desenvolvido sobre o conjunto de dados SemClinBr. A partir desse problema, estabelece-se a seguinte pergunta de pesquisa:

\begin{quote}
\emph{Como desenvolver o RECLin-PT para a extração de relações clínicas em português
no corpus SemClinBr, com ênfase na identificação de relações de negação?}
\end{quote}

Responder a essa pergunta envolve decisões para as quais a literatura não oferece
resposta pronta, como a escolha do modelo de linguagem pré-treinado sobre o qual o
RECLin-PT é construído e a forma de organizar a decisão de classificação. Neste
trabalho, essas decisões são investigadas experimentalmente, como parte do próprio
desenvolvimento, para definir uma configuração adequada ao modelo. O desafio é agravado pelo forte desbalanceamento entre as classes de relação e pelo fato de que o SemClinBr contém apenas anotações de relações positivas, o que exige a construção controlada de exemplos de ausência de relação.

\section{Motivação}
\label{sec:motivacao}

A negação clínica é um fenômeno linguístico de alto impacto, o que justifica um
modelo com ênfase nessa relação. Em português, ela é fortemente lexicalizada (``não'',
``sem'', ``nega'', ``ausência de''), e há a expectativa difundida de que modelos
pré-treinados sobre texto do próprio domínio clínico capturem melhor essas
construções do que modelos de domínio geral. Esse ganho, contudo, não é automático. Em
inglês, \citeonline{lin_2020_negacao} reportam que, na detecção de negação clínica, o
BERT sem as extensões de adaptação de domínio avaliadas obteve, em média, o melhor
desempenho. Para o português, não foram encontradas avaliações controladas que permitam isolar o efeito do pré-treinamento de domínio na extração de relações clínicas. A comparação direta
entre modelos clínicos e de domínio geral sobre o SemClinBr foi conduzida para o
reconhecimento de entidades \cite{almeida_2026_benchmarkner}, e não para a etapa de
relações (Capítulo~\ref{cap:relacionados}). Desenvolver o RECLin-PT exige, portanto,
examinar essa escolha em vez de pressupô-la.

A escolha do modelo pré-treinado também não é independente da forma como a decisão de
classificação é organizada, e alternativas distintas só podem ser comparadas de forma
justa se passarem pelos mesmos dados e pelo mesmo protocolo de seleção e de
avaliação. Por isso, o desenvolvimento do RECLin-PT apoia-se em uma cadeia
experimental reprodutível, que também responde a críticas recorrentes à literatura
de PLN biomédico, como o compartilhamento restrito de modelos e o detalhamento
insuficiente do ajuste fino\footnote{Ajuste fino (\foreignlanguage{english}{fine-tuning})
é o retreinamento de um modelo pré-treinado sobre os dados de uma tarefa específica,
partindo dos pesos já aprendidos em vez de uma inicialização aleatória
\cite{devlin_2019_bert}.} \cite{navarro_2023_revisao,shimaoka_2026_revisao}.

\section{Objetivos}

\subsection{Objetivo geral}

Desenvolver o RECLin-PT, um modelo de extração de relações clínicas em português
aplicado ao corpus SemClinBr, com ênfase na identificação de relações de negação
(\texttt{negation\_of}), investigando alternativas de modelo pré-treinado e de
estratégia de modelagem para subsidiar a definição de sua configuração.

\subsection{Objetivos específicos}

\begin{enumerate}
  \item Preparar os dados do corpus SemClinBr para a tarefa de extração de relações,
        considerando o alinhamento entre os textos e as anotações e a organização de
        conjuntos experimentais de treino, desenvolvimento e teste.
  \item Caracterizar as relações clínicas presentes no corpus e estabelecer
        procedimentos reprodutíveis para a geração de pares candidatos de entidades.
  \item Desenvolver o RECLin-PT, investigando representações de pares de entidades e
        estratégias de modelagem para a classificação de relações clínicas.
  \item Avaliar experimentalmente as configurações consideradas, utilizando
        diferentes modelos pré-treinados e métricas adequadas a uma tarefa com
        classes desbalanceadas.
  \item Definir uma configuração de referência para o RECLin-PT e estabelecer os
        procedimentos para o congelamento de sua configuração final e para a
        avaliação única dessa configuração no conjunto de teste.
\end{enumerate}
